
\subsubsection{Overview}

Given an attribution method A, a trained model M, and a set of contrastive probe pairs P, the method's sensitivity to each influence mode is measured by computing attribution scores on probes designed to isolate that mode. The resulting mode profile is a vector p\_A = [s\_mem, s\_transfer, s\_spurious] capturing the method's relative sensitivity to memorization, feature transfer, and spurious association.

\subsubsection{Influence Modes}

Three modes of training data influence are defined:

\textbf{Memorization:} Training examples influence test predictions through direct memorization---the model learns the specific input-output mapping rather than generalizable features. Attribution methods sensitive to memorization assign high influence to near-duplicate or highly similar training examples.

\textbf{Feature Transfer:} Training examples contribute features that generalize to test examples in the same category. Attribution methods sensitive to feature transfer assign influence based on shared high-level representations rather than surface similarity.

\textbf{Spurious Association:} Training examples share spurious correlations (e.g., background features, co-occurring attributes) with test examples. Attribution methods sensitive to spurious association may attribute influence to examples sharing these shortcuts.

\subsubsection{Contrastive Probe Construction}

For each mode m $\in$ {mem, transfer, spurious}, contrastive probe pairs (t, $s^+, s^-)$ are constructed where t is a test point, $s^+$ is a training point exhibiting mode m relationship with t, and $s^-$ is a matched control lacking this relationship.

\textbf{Memorization probes:} $s^+$ selected as training examples with high pixel-level similarity to t (near-duplicates). $s^-$ selected from the same class with low surface similarity.

\textbf{Feature transfer probes:} $s^+$ selected from the same semantic class as t but with different surface features. $s^-$ selected from a different class with similar surface statistics.

\textbf{Spurious probes:} $s^+$ selected sharing a spurious correlation with t. $s^-$ selected from the same class without the spurious correlation.

Mode sensitivity is computed as:

$s_m = (1/|P_m|) \Sigma sign(A(s^+; t) - A(s^-; t))$

where A(s; t) is the attribution score for training point s given test point t.

\subsubsection{Dissociation Analysis}

To test whether methods produce dissociable fingerprints, each method is run with multiple random seeds and ANOVA analysis is applied:

\textbf{Between-group variance (inter-method):} Variance of method means around the grand mean, capturing how different methods' average profiles are.

\textbf{Within-group variance (intra-method):} Variance of individual runs around each method's mean, capturing random variation across seeds.

The F-ratio F = MS\_between / MS\_within quantifies dissociation. An F-ratio significantly greater than 1 indicates that method identity explains more variance than random seed---i.e., methods produce systematic fingerprints.

Thresholds are set based on standard statistical practice: F > 4.0 for meaningful dissociation (corresponds to p < 0.05 with the experimental degrees of freedom), and Cohen's d > 0.5 for medium-to-large effect size in pairwise comparisons.

\subsubsection{Cross-Architecture Transfer}

To test whether fingerprints are method-intrinsic rather than model-specific, mode profiles are computed for each method across multiple architectures (ResNet-18, ViT-Small, ConvNeXt-Tiny) and Pearson correlation is measured between profiles of the same method on different architectures.

High correlation (r > 0.7) would indicate that a method's fingerprint transfers across architectures.

\subsubsection{Attribution Methods}

Three methods spanning different computational approaches are evaluated:

\textbf{TRAK} \citep{park2023trak}: Computes influence via random projection of gradients, preserving gradient direction while reducing dimensionality.

\textbf{TracIn} \citep{pruthi2020tracin}: Sums gradient dot products weighted by learning rate across checkpoints. Captures temporal patterns in how influence accumulates during training.

\textbf{Kronfluence} \citep{grosse2023kronfluence}: Inverts Fisher information matrix using K-FAC approximation. Emphasizes curvature structure in the loss landscape.
